\documentclass[11pt]{article}
\usepackage[margin=1in]{geometry}
\usepackage{amsmath,amssymb,amsthm}
\usepackage{booktabs}
\usepackage{hyperref}
\usepackage[numbers,sort&compress]{natbib}
\hypersetup{hidelinks}
\emergencystretch=3em
% Long Lean identifiers: allow \path to break after underscores.
\expandafter\def\expandafter\UrlBreaks\expandafter{\UrlBreaks\do\_}

\newtheorem{theorem}{Theorem}
\newtheorem{proposition}{Proposition}
\newtheorem{corollary}{Corollary}
\newtheorem{lemma}{Lemma}
\newtheorem{definition}{Definition}
\newtheorem{remark}{Remark}
\newtheorem{conjecture}{Conjecture}
\newtheorem{observation}{Observation}
\newtheorem{obligation}{Proof Obligation}

% A clearly-marked aside for readers who DO know the topology.  Everyone
% else can skip these blocks without losing the thread.
\newenvironment{cohomaside}
  {\par\medskip\noindent\begingroup\small
   \textbf{[For the cohomologist.]}\ }
  {\endgroup\par\medskip}

\newcommand{\Hyp}{\Xi}
\newcommand{\feas}{\mathrm{feas}}
\newcommand{\Sec}{\mathcal S}
\newcommand{\rectrho}{\rho}
\newcommand{\obs}{\mathrm{obs}}
\newcommand{\Hech}{\check H}

\title{Local Pieces, Global Solutions:\\
A Machine-Checked Sheaf-Cohomological Account of\\
Set-Theoretic Estimation\\
\large From Combettes' feasibility sets and Carlson's cryptanalytic
STE to a mechanized gluing theorem,\\ with an empirical corpus
instantiation --- written to be readable without prior cohomology}
\author{Project \texttt{thejeb}}
\date{July 28, 2026}

\begin{document}
\maketitle

\begin{abstract}
Set-theoretic estimation (STE) answers an estimation problem by
intersecting property sets: an estimate is acceptable exactly when it
is consistent with every piece of information
\citep{combettes1993}.  When the hypothesis space is a product over
$N$ variables, the feasibility set lives behind an exponential wall,
and the standard escape is \emph{locality}: keep small tables over
small, overlapping sets of variables.  Every local scheme lives or
dies on one question --- \emph{when does locally consistent partial
information determine a global solution?}  The mathematics of that
question is sheaf theory and \v{C}ech cohomology
\citep{abramsky2011sheaf,abramsky2012cohomology}, and this paper
develops it for STE in machine-checked form (Lean~4 / Mathlib
\citep{moura2021,mathlib2020}) while assuming no topology of the
reader: every term --- cover, section, compatible family, gluing,
obstruction --- is defined from STE primitives, with its
constraint-satisfaction meaning stated in one sentence beside it.  The
main theorem, mechanized sorry-free in the module
\path{Ste.SupportCover}, is the \emph{support--cover gluing theorem}:
over any cover of the variables, every compatible family of local
sections glues to a global solution \emph{iff} the constraint is
generated by pieces whose scopes fit inside the cover --- so gluing
decides representability over given scopes, and the entire
computational difficulty of STE relocates into \emph{computing local
sections}, with gluing itself free.  Around this centerpiece we prove
exact obstruction counts on coupling constraints, rectangle-cover and
fooling-set bounds with machine-checked incomparability results, and a
reading of the bounded-treewidth algorithmic stack (bucket
elimination, junction trees, adaptive consistency) as the theory of
when local sections are cheap.  We connect the picture to Carlson's
cryptanalytic instantiation of STE \citep{carlson2012}, whose
mechanized counting and reduction theorems we reinterpret sheaf-wise
--- honestly separating what maps, what is dispatched to in-flight
mechanization, and what is speculation.  Finally we report a
constructive proof-of-work: a verified solver whose pairwise
consistency check ran on three historical text corpora, together with
the mechanized bridge theorems showing that for frame-structured
corpora this pairwise check is a complete feasibility test --- the
cohomological theory explains \emph{why} the cheap check suffices.
Throughout, a strict bright line separates machine-checked results
(each cites its Lean identifier) from conjecture, and every heavy
cochain construction is quarantined in a clearly marked appendix.
\end{abstract}

\section{Introduction}
\label{sec:intro}

This paper tells one story three times, at increasing precision: in
plain constraint-satisfaction language, in the vocabulary of sheaves
built up from STE primitives, and in the named theorems of a Lean~4
mechanization.  The story: an STE problem intersects property sets;
intersections over product spaces hide an exponential wall; local
reasoning over small variable scopes is the escape; local reasoning is
sound and complete exactly when local pieces \emph{glue}; and the
gluing question has an exact, machine-checked answer --- it succeeds
precisely when the constraint is cut out by pieces that fit the
chosen scopes.  Everything else in the paper --- obstruction counts,
representation bounds, the algorithmic stack, the cryptanalytic and
corpus instantiations --- is that one theorem seen from different
distances.

\paragraph{Contributions.}
\begin{enumerate}
\item An \emph{accessible on-ramp} to the sheaf view of STE: every
  topological word is defined from STE primitives before use
  (Sections~\ref{sec:background}--\ref{sec:dictionary}), with worked
  examples small enough to check by hand.
\item The \emph{support--cover gluing theorem}
  (Section~\ref{sec:main}), mechanized sorry-free
  (\path{Ste.SupportCover}): the \v{C}ech obstruction of a cover
  vanishes iff the constraint admits a generating family refining the
  cover.  Corollaries: gluing decides representability over given
  scopes; the obstruction is a certificate of non-decomposability;
  and (marked as sketches) a ``gluing width'' invariant equal to
  generating arity, provably distinct from treewidth.
\item Machine-checked \emph{quantitative} results
  (Section~\ref{sec:main}): exact obstruction counts on coupling
  constraints, rectangle-cover and fooling-set bounds, and
  incomparability theorems refuting the slogan ``obstruction size $=$
  representation blow-up.''
\item An honest \emph{Carlson connection}
  (Section~\ref{sec:carlsonconn}): which parts of the dissertation's
  STE cryptanalysis \citep{carlson2012} map onto the sheaf picture
  (with mechanized witnesses), which parts are dispatched to
  in-flight mechanization, and which do not map.
\item An \emph{empirical instantiation} (Section~\ref{sec:empirical}):
  a verified solver run on three historical corpora, with the
  mechanized equivalence between its pairwise check and genuine
  feasibility --- and the cohomological explanation of why pairwise
  suffices for this constraint class.
\end{enumerate}

\paragraph{The bright line.}  One convention is enforced throughout
and is non-negotiable.  Every machine-checked claim cites its Lean
identifier in \texttt{this font}; the named theorem exists, sorry-free,
in the project library under continuous integration.  Everything not
so cited is explicitly labelled: \emph{conjecture}, \emph{sketch},
\emph{dispatched} (an exact statement handed to an in-flight
mechanization effort, not yet part of the trusted base), or
\emph{speculation}.  One refinement: a small number of results were
mechanized, during the writing of this paper, on named research
branches not yet merged into the main library; these are cited with
their branch and commit stated inline, after independent
verification that the branch file is sorry-free and compiles, and
they are never silently mixed with main-library citations.  Adding
intuition never relaxes this line; intuition sits beside the
citations, not in place of them.

\paragraph{How to read this paper.}  A reader who knows STE or
constraint satisfaction but no algebraic topology can read
straight through: Sections~\ref{sec:ste}--\ref{sec:why} need no new
vocabulary, Section~\ref{sec:background} builds the five words the
rest uses, and no commutative diagram appears outside the appendix.
Asides marked \textbf{[For the cohomologist]} and
Appendix~\ref{app:coefficients} carry the cochain-level machinery and
can be skipped without losing the thread.

\section{Set-theoretic estimation: foundations}
\label{sec:ste}

STE, as founded by \citet{combettes1993} (with set-membership
antecedents in \citealp{schweppe1968}), replaces optimization by
\emph{feasibility}.  An estimation problem is specified by a solution
space $\Hyp$, an index set $I$ of pieces of information, and for each
$i \in I$ a \emph{property set} $S_i \subseteq \Hyp$: the estimates
consistent with that piece of information.  The answer object is the
\emph{feasibility set}
\[
  T \;=\; \bigcap_{i \in I} S_i,
\]
and a set-theoretic estimate is any point of $T$.  The mechanization
of this core is \path{Ste.Basic}: \path{feasibilitySet},
\path{mem_feasibilitySet}, the refinement law
\path{feasibilitySet_subset}, partial enforcement
\path{partialFeasibilitySet} with its monotonicity
\path{partialFeasibilitySet_antitone} (more information never enlarges
the feasible set), and Combettes' taxonomy: a formulation is
\emph{consistent} if $T \ne \emptyset$, \emph{fair} if the true object
lies in $T$ (\path{Fair}, whence consistency:
\path{feasibilitySet_nonempty_of_fair}), and \emph{ideal} if $T$ is
exactly the true object (\path{Ideal}, \path{Ideal.eq_of_mem}).

Throughout this paper the hypothesis space is a product: $\Hyp =
\prod_{v \in V} A_v$, a hypothesis assigning a value to each of $N$
variables, each property set constraining some of them.  This is the
constraint-satisfaction shape of STE, and it hides a wall.  The
feasible set lives in a space of size exponential in $N$, and the
project has machine-checked that no representation cleverness can
help in general: crisp property sets realize \emph{any} of the $2^N$
extensional states (\path{card_allExactStates},
\path{arbitrary_subset_as_single_property}, in
\path{Ste.DynamicFrame.FiniteSolver}).  Any solver that manipulates
$T$ extensionally pays $2^N$.  The rest of the paper is about the
principled way around: reason \emph{locally}, and understand exactly
when local reasoning is complete.

\section{Carlson and the topological turn}
\label{sec:carlson}

Between Combettes' analytic setting (metric and Hilbert spaces,
projections onto convex sets) and this paper's combinatorial one sits
Carlson's dissertation \citep{carlson2012}, which recast STE over
finite, discrete \emph{topological} spaces and applied it to
cryptanalysis.  The move matters for us twice over: it is the
project's motivating instantiation of STE outside signal processing,
and it is the first place in this lineage where the \emph{topological}
character of feasibility --- closed sets, compactness, intersection
structure --- is taken seriously as the load-bearing mathematics.
This paper's sheaf layer is the continuation of exactly that turn:
from the topology of the feasible set to the topology of \emph{where
information lives}.

\subsection{Decryption as STE}

Carlson's recast, mechanized in \path{Ste.Carlson.Cipher}: the
solution space is a finite key space; each observed ciphertext $c$
contributes the property set of keys under which $c$ decrypts to a
meaningful message; decryption succeeds when the feasibility set
collapses to the true key --- Shannon's unicity point
\citep{shannon1949} in set-theoretic form.  Machine-checked: the
cipher-system skeleton (\path{CipherSystem}, unambiguous decryption
\path{decrypt_unique}), the induced property sets
(\path{keyPropertySet}) and feasible keys (\path{feasibleKeys},
\path{mem_feasibleKeys}), fairness of genuine traffic --- if the
observations really are encryptions under a true key, that key is
feasible (\path{fair_of_genuine},
\path{feasibleKeys_nonempty_of_genuine}) --- unicity as ideality
(\path{Unicity}, \path{Unicity.eq_of_feasible}), and observation
monotonicity: more intercepted traffic can only narrow the key search
(\path{feasibleKeys_antitone}, a special case of
\path{partialFeasibilitySet_antitone}).

The topological well-posedness facts Carlson's setting needs are
mechanized in \path{Ste.Topology}: closed property sets give a closed
feasibility set (\path{isClosed_feasibilitySet}; discrete case
\path{isClosed_feasibilitySet_of_discrete}), fairness gives the
finite intersection property (\path{FiniteConsistency.of_fair}), and
compactness turns finite consistency into consistency outright
(\path{feasibilitySet_nonempty_of_compact}); ideal formulations have
closed feasible sets in any $T_1$ space
(\path{Ideal.isClosed_feasibilitySet}).

\subsection{Counting and reduction, machine-checked}

Two quantitative pillars of the dissertation are mechanized exactly.

\emph{Residual key counting} (\path{Ste.Carlson.Counting}).  For a
substitution cipher over an alphabet $A$ --- keys are permutations of
$A$ --- an intercepted message constrains the key only on the set $T$
of symbols that actually occur.  Carlson's Lemma~4.1 counts the keys
equivalent to the true key: exactly $(|A| - |T|)!$ of them
(\path{card_consistent_keys}, via \path{card_perm_fixing}; the key
space itself has size $|A|!$, \path{card_substitution_keys}).  This is
the size of the STE feasible key set after one observation --- the
residual ambiguity.

\emph{Cipher reduction} (\path{Ste.Carlson.Reduction}).  On $n$-bit
blocks, every permutation-cipher key (a permutation of bit positions)
is realized as a substitution-cipher key (a permutation of the value
space) by the induced map \path{coperm}
(\path{permutation_reduces_to_substitution}), injectively
(\path{coperm_injective}); but for $n \ge 2$ there are strictly more
substitution keys than permutation keys
(\path{card_permutation_lt_card_substitution}), so no faithful
embedding is onto (\path{substitution_not_reducible_to_permutation})
--- Carlson's Lemma~5.1/Corollary~5.2 --- and
permutation--substitution--permutation compositions collapse to a
single substitution key (\path{psp_reduces_to_substitution},
Corollary~5.3; the same reduction is developed in
\citealp{carlson2022equivalence,ghosh2021isomorphic}).

Section~\ref{sec:carlsonconn} returns to Carlson with the sheaf
vocabulary in hand and asks, honestly, which of these constructions
are secretly about local sections and gluing.

\section{Why sheaves and cohomology are a useful play}
\label{sec:why}

The escape from the $2^N$ wall that every practical system reaches
for is \emph{locality}: instead of the monolithic $T$, keep small
tables about small, overlapping sets of variables --- per-variable
domains, arc-consistent pair tables, junction-tree bags --- and
reason with those.  Every such scheme lives or dies on one question:

\begin{quote}
\emph{If I know everything the problem permits locally --- on each
small set of variables --- and my local answers never contradict each
other, do they determine a genuine global solution?}
\end{quote}

Sometimes yes, sometimes no.  When yes, local reasoning is
\emph{complete}: the cheap, small tables are a faithful proxy for the
exponentially large $T$.  When no, the local tables admit
\emph{phantoms} --- mutually consistent local answers realized by no
global solution --- and any margins-only procedure wastes work on
them, or worse, reports them.  The mathematical bookkeeping for
exactly this local-versus-global question is called sheaf theory, and
the device that counts the failures is called \v{C}ech cohomology.
The names do not matter; the payoff does:

\begin{itemize}
\item \textbf{An exact criterion} for when local reasoning over given
  scopes is complete (the support--cover theorem,
  Section~\ref{sec:main}) --- machine-checked, with no structural
  side conditions on how the scopes overlap.
\item \textbf{Certificates}: a single stuck family of local answers
  certifies that \emph{no} decomposition of the constraint respects
  those scopes; a fooling set certifies that no small
  union-of-rectangles representation exists
  (Section~\ref{sec:quant-sub}).
\item \textbf{A precise location for hardness}: the theory proves
  gluing is never the expensive step --- computing genuine local
  tables is --- which turns the bounded-width algorithmic stack into
  a completeness theorem for local reasoning
  (Section~\ref{sec:algorithms-sub}).
\item \textbf{A no-free-lunch theorem}: any invariant computed from
  the shape of the scopes alone, ignoring the tables, is provably
  useless (machine-checked; Appendix~\ref{app:constant}) --- which
  tells you where \emph{not} to look for speedups.
\end{itemize}

Where can the local-to-global question even fail?  The development
contains three distinct axes of locality, and it is worth fixing at
the outset that the question is only genuinely hard on one of them.

\begin{enumerate}
\item \textbf{Splitting the constraint list} (index set $I$).
  Partition or arbitrarily cover the constraints, intersect within
  groups, intersect the results: you always recover exactly $T$.  The
  map $J \mapsto \feas(J) = \bigcap_{i \in J} S_i$ sends unions of
  constraint sets to intersections of partial feasible sets
  (\path{partialFeasibilitySet_iUnion}), so any cover of the
  constraint index recovers the global feasible set
  (\path{feasibilitySet_eq_iInter_cover}, in \path{Ste.Sheaf}).
  \emph{Local-to-global holds here for free}; nothing can obstruct
  it.  (Under the formal-concept reading of \path{Ste.HyperFrame},
  $\feas$ is a lower polar of the satisfaction context
  \citep{ganter1999fca} (\path{partialFeasibilitySet_eq_lowerPolar}),
  and polars always turn joins into meets.)
\item \textbf{Splitting the variable set} (scopes $W \subseteq V$).
  Keep, for each small scope, the partial assignments the problem
  permits.  Here local-to-global genuinely fails, the failure is
  precisely variable \emph{coupling}, and this is where the paper
  lives.
\item \textbf{Splitting the evidence corpus} (document sets $D$, in
  \path{Ste.DynamicFrame}).  Feasibility over $D \cup E$ splits
  exactly: feasible on the union iff feasible on each piece
  \emph{and} the finite, explicitly computable set of constraints
  active only on the union is satisfied (\path{crossObstruction},
  \path{mem_feasible_union_iff}, in
  \path{Ste.DynamicFrame.PresheafGluing}); a common extension of
  compatible local worlds exists iff that correction set is empty,
  and is then unique (\path{globalExtension_nonempty_iff},
  \path{globalExtension_subsingleton}, \path{glue}).  A
  local-to-global theorem \emph{with an exact error term}, already
  closed.
\end{enumerate}

The slogan the mechanization supports: \emph{the constraint side
glues for free, the corpus side glues with an exact and computable
correction term, and the variable side is where the real
local-versus-global phenomenon occurs}.  From here on, ``local''
means local in the variables.

\section{Related work}
\label{sec:related}

\paragraph{Sheaf-theoretic contextuality.}  The program closest to
this paper is Abramsky--Brandenburger's sheaf-theoretic treatment of
non-locality and contextuality \citep{abramsky2011sheaf}, its
cohomological refinement \citep{abramsky2012cohomology}, the logical
reading via paradox \citep{abramsky2015paradox}, contextual semantics
across disciplines \citep{abramsky2014contextualsemantics}, the
quantitative contextual fraction \citep{abramsky2017contextualfraction},
and the studies of when cohomological invariants are complete or have
false negatives
\citep{caru2017cohomology,caru2018complete,okay2017topological,aasnaess2022comparing}.
The present work differs in two ways: the section presheaf of STE
consists of \emph{globally extendable} partial assignments (their
empirical models are arbitrary compatible data), which collapses the
per-section obstruction (\path{amb_extension_always},
Appendix~\ref{app:twisted}) and moves the interesting content to
family-level gluing; and every result here claimed as proven is
machine-checked in Lean~4 \citep{moura2021,mathlib2020}.

\paragraph{Structural constraint satisfaction.}  The
algorithmic layer this paper reinterprets is classical: consistency
enforcement \citep{mackworth1977consistency,freuder1982backtrack},
constraint networks and directional/adaptive consistency
\citep{dechterpearl1987network}, tree clustering
\citep{dechterpearl1989tree}, bucket elimination
\citep{dechter1999bucket,dechter2003constraint}, with the width
theory of \citet{bertele1972nonserial},
\citet{robertson1986graphminorsII}, \citet{arnborg1987complexity},
\citet{bodlaender1998arboretum}, and the homomorphism-side
tractability characterization of \citet{grohe2007complexity}.  The
paper's contribution to this literature is not a new algorithm but a
mechanized exact boundary: which covers admit complete local
reasoning at all, independent of acyclicity hypotheses.

\paragraph{Representation lower bounds.}  The rectangle-cover
number and fooling sets are imported from communication complexity
\citep{kushilevitz1997communication,aho1983notions}; the extension to
LP extension complexity
\citep{yannakakis1991expressing,fiorini2015exponential} is flagged as
outlook.

\paragraph{STE lineage.}  Foundations \citep{combettes1993},
set-membership estimation \citep{schweppe1968}, and the cryptanalytic
line \citep{shannon1949,carlson2012,ghosh2021isomorphic,
carlson2022equivalence}.

\section{Background: cohomology, built from STE primitives}
\label{sec:background}

This section defines every piece of vocabulary the paper uses.  Each
definition is an ordinary constraint-satisfaction object; the
topological name is recorded so that the Lean files and the
literature can be read, but nothing later depends on knowing why
topologists chose it.

\subsection{The five words}
\label{sec:fivewords}

\paragraph{Scope, and cover.}  A \emph{scope} (the Lean files say
\emph{context}) is a subset $W \subseteq V$ of the variables.  A
\emph{cover} is a choice of scopes $U : J \to \mathcal P(V)$ such
that every variable appears in at least one ($\forall v\,\exists j,\
v \in U_j$).  That is all a cover is: \textbf{a choice of overlapping
variable-scopes to reason locally over}.  Three examples recur: the
\emph{singleton cover} $\{\{v\} : v \in V\}$ (pure per-variable
reasoning); the \emph{bag cover} given by the bags of a junction tree
or elimination order; the \emph{trivial cover} $\{V\}$ (no locality
at all).

\paragraph{Local section.}  For a constraint $T \subseteq \prod_v
A_v$ and a scope $W$, a \emph{local section of $T$ over $W$} is
\textbf{a partial assignment on $W$ that extends to a full solution
of $T$} --- the restriction to $W$ of some member of $T$.  The set of
them, $\Sec_T(W) = {}$\path{localSections}$\,T\,W$, is the image of
$T$ under coordinate restriction (\path{Ste.VariablePresheaf}).  Note
the convention carefully: a local section is not merely ``locally
plausible''; global extendability is baked into the definition.
(Partial assignments that are only locally consistent --- what a
propagation algorithm computes --- are a weaker object;
Section~\ref{sec:pseudo} names them \emph{pseudo-sections}, and the
distinction carries much of the story.)  Restriction behaves as it
should --- restricting in two steps equals restricting in one
(\path{restrict_restrict}, \path{localSections_restrict₂}).  The
topologist's word for ``a scope-indexed family of section sets with
well-behaved restriction'' is a \emph{presheaf}; that bookkeeping is
all the word means here.

\paragraph{Compatible family.}  Fix a cover $U$.  A \emph{compatible
family} is a choice of one local section $s_j \in \Sec_T(U_j)$ per
scope such that any two agree on the overlap of their scopes
(\path{CompatibleFamily}, in \path{Ste.CechCover}): \textbf{local
answers that never contradict each other on shared variables}.  This
is exactly the state of knowledge of a solver that has computed
perfect local tables and checked them pairwise.

\paragraph{Gluing.}  A compatible family \emph{glues}
(\path{GluesCover}) when a single global solution $f \in T$ restricts
to $s_j$ on each $U_j$: \textbf{the local pieces can be stitched into
one global solution}.  When every compatible family over $U$ glues,
the \emph{sheaf condition} holds for $U$ (\path{CechVanishesCover});
in topological notation ``$\Hech^1$ vanishes for the cover,'' and its
STE meaning is exactly \textbf{``every locally consistent choice is
globally realizable.''}

\paragraph{Obstruction.}  When the sheaf condition fails, the failure
can be counted:
\[
  \obs \;=\; \#\{\text{compatible families}\} \;-\;
             \#\{\text{families that glue}\},
\]
\textbf{the number of ways local answers can be mutually consistent
yet stitchable into no global solution}.  The singleton-cover count is
\path{cechObstruction}; the cover-level count is
\path{coverCechObstruction}, and it vanishes whenever the sheaf
condition holds (\path{coverCechObstruction_eq_zero}).  Each stuck
family is a \emph{phantom}: a profile of partial assignments that
every margins-only procedure must entertain although no solution
produces it.

\medskip
\noindent Table~\ref{tab:dictionary} collects the dictionary.

\begin{table}[ht]
\centering
\footnotesize
\begin{tabular}{@{}p{9.5em}p{15em}p{12.5em}@{}}
\toprule
topological word & STE meaning & Lean name \\
\midrule
context / open set & scope: a subset $W \subseteq V$ of variables & --- \\
cover & choice of overlapping scopes hitting every variable &
  (hypothesis of \path{CechVanishesCover}) \\
local section over $W$ & partial assignment on $W$ extending to a
  solution & \path{localSections} \\
presheaf & scope-indexed section sets, consistent restriction &
  \path{restrict_restrict} \\
compatible family & local answers agreeing on shared variables &
  \path{CompatibleFamily} \\
gluing & stitching the pieces into one global solution &
  \path{GluesCover} \\
sheaf condition ($\Hech^1 = 0$) & every compatible family glues &
  \path{CechVanishesCover} \\
\v{C}ech obstruction & count of compatible-but-unstitchable families &
  \path{cechObstruction}, \path{coverCechObstruction} \\
\bottomrule
\end{tabular}
\caption{The five-word dictionary.  Everything in this paper is one
of these objects.}
\label{tab:dictionary}
\end{table}

\paragraph{Why would a compatible family ever fail to glue?}  Each
piece individually extends to a global solution --- that is the
definition of a local section.  The catch: each piece may extend to a
\emph{different} solution, with no single solution extending all of
them at once.  That one sentence is the entire phenomenon; the next
subsection exhibits it in four rows of a truth table.

\subsection{The smallest example, worked by hand}
\label{sec:example}

All statements in this subsection and the next are proved in Lean,
sorry-free (the singleton-cover computations live in
\path{Ste.CechObstruction}); theorem names are cited exactly.

\paragraph{Two Booleans that must agree.}  Take two Boolean variables
$x, y$ and the single constraint ``$x = y$'':
\[
  \Delta \;=\; \{f : \{x,y\} \to \mathsf{Bool} \mid f(x) = f(y)\}
\]
(\path{diagonal}, in \path{Ste.Sheaf}).  It has exactly two
solutions, $(\mathsf{ff},\mathsf{ff})$ and
$(\mathsf{tt},\mathsf{tt})$ (\path{diagonal_encard}).  Reason locally
over the singleton cover: one scope $\{x\}$, one scope $\{y\}$.

\emph{Local sections.}  Over $\{x\}$: does $x = \mathsf{ff}$ extend
to a solution?  Yes ($(\mathsf{ff},\mathsf{ff})$).  Does $x =
\mathsf{tt}$?  Yes ($(\mathsf{tt},\mathsf{tt})$).  So both values of
$x$ are local sections, and symmetrically for $y$
(\path{contextSections_diagonal}): every partial view is fully
uninformative --- each variable, seen alone, can still be anything.

\emph{Compatible families.}  The two scopes do not overlap
(\path{singletonCover_overlap_empty}), so there is nothing to agree
on: compatibility is automatic, and a compatible family is just a
choice of one value per variable --- four families in all
(\path{compatibleFamilies},
\path{compatibleFamilies_diagonal_encard}):
\[
(\mathsf{ff},\mathsf{ff}) \quad (\mathsf{ff},\mathsf{tt}) \quad
(\mathsf{tt},\mathsf{ff}) \quad (\mathsf{tt},\mathsf{tt}).
\]

\emph{Which glue?}  A family glues iff, read as a full assignment, it
is a solution (\path{glues_iff_mem}, \path{gluedFamilies_eq}).  Two
do (\path{gluedFamilies_diagonal_encard}).  The mixed family $x
\mapsto \mathsf{ff},\ y \mapsto \mathsf{tt}$ is the interesting one:
\textbf{each half extends to a solution, but to different solutions},
and no single solution extends both.  It is compatible and stuck,
machine-checked as \path{diagonal_mixed_compatible_not_glues}.  The
obstruction count is
\[
  \obs(\Delta) \;=\; 4 - 2 \;=\; 2
\]
(\path{diagonal_cechObstruction}).  Two phantoms: the two mixed
families.  This is what a nonzero \v{C}ech obstruction \emph{is} ---
nothing more exotic than the mixed rows of this truth table.

\paragraph{The exact boundary: gluing iff no coupling.}  Call $T$
\emph{rectangular} if it is a product of per-variable sets, $T =
\prod_v P_v$ (\path{IsRectangle}) --- the constraint imposes no
coupling, only per-variable filters.  Rectangular problems are the
fragment where per-variable reasoning is complete: the feasible set
of a rectangular family is itself a rectangle
(\path{rectangular_feasibilitySet}) of cardinality $\prod_v
|{\bigcap_i P_{i,v}}|$ (\path{rectangular_encard}), representable in
$\sum_v$ space, and its singleton-cover obstruction is zero
(\path{rectangular_cechObstruction}).  The diagonal is provably not
rectangular (\path{diagonal_not_rectangular}); worse, the smallest
rectangle containing it is the whole four-point space
(\path{rectangle_superset_diagonal_encard}) --- the tightest
per-variable view of a coupling carries \emph{zero} information.  The
singleton-cover story is closed as an iff
(\path{cechVanishes_iff_rectangular}, in
\path{Ste.CouplingLowerBound}):
\[
  \text{every compatible family glues}
  \;\iff\;
  T \text{ is a rectangle.}
\]
In plain terms: \textbf{per-variable local reasoning is exactly as
strong as the absence of coupling}.  At this cover the obstruction
does not merely notice coupling; it characterizes it.

\paragraph{Scaling: $n$ variables that must all agree.}  The phantom
count is not an artifact of small examples.  The $n$-fold coupling
$\mathrm{allEqual}(n, \alpha)$ --- all $n$ coordinates equal, over an
alphabet $\alpha$ (\path{allEqual}) --- has every per-variable margin
full (\path{contextSections_allEqual}), exactly as the diagonal did:
locally, anything goes.  So \emph{all} $|\alpha|^n$ families are
compatible (\path{compatibleFamilies_allEqual_encard}), while only
the $|\alpha|$ constant assignments are solutions
(\path{allEqual_encard}).  The obstruction is exactly $|\alpha|^n -
|\alpha|$ (\path{cechObstruction_allEqual}), hence at least $2^n - 2$
once $|\alpha| \ge 2$ (\path{cechObstruction_allEqual_ge}):
exponentially many phantoms while the true feasible set keeps
constant size.  For $n \ge 2$ the constraint is genuinely
non-rectangular (\path{allEqual_not_rectangular},
\path{allEqual_not_cechVanishes}).  A margins-only pipeline on this
instance does exponential work distinguishing phantoms from solutions
--- the $2^N$ wall of Section~\ref{sec:ste}, seen from the local
side.

\section{The dictionary, precisely}
\label{sec:dictionary}

For reference, this section restates the correspondence at full
precision; it adds no new ideas beyond
Section~\ref{sec:background}.

\begin{definition}[The section presheaf; mechanized,
\path{Ste.VariablePresheaf}]
For $T \subseteq \prod_{v \in V} A_v$ and $W \subseteq V$, the local
sections are
$\Sec_T(W) = \{\, f|_W \mid f \in T \,\}$
(\path{localSections}, the image of $T$ under $f \mapsto (v : W)
\mapsto f\,v$).  For $W' \subseteq W$ the restriction map
$\Sec_T(W) \to \Sec_T(W')$ is coordinate restriction; functoriality
is \path{restrict_restrict} and \path{localSections_restrict₂}.  This
makes $W \mapsto \Sec_T(W)$ a contravariant functor on scopes --- a
presheaf of sets.
\end{definition}

\begin{definition}[Compatibility, gluing, vanishing; mechanized,
\path{Ste.CechCover}]
For a cover $U : J \to \mathcal P(V)$: a family $(s_j)_{j \in J}$
with $s_j \in \Sec_T(U_j)$ is \emph{compatible}
(\path{CompatibleFamily}) if $s_j$ and $s_k$ agree on $U_j \cap U_k$
for all $j,k$; it \emph{glues} (\path{GluesCover}) if some $f \in T$
has $f|_{U_j} = s_j$ for all $j$; the constraint satisfies the sheaf
condition for $U$ (\path{CechVanishesCover}) if every compatible
family glues.  Glues along a cover are unique
(\path{glueCover_unique}: two solutions agreeing on every scope of a
cover agree everywhere), so the sheaf condition is genuinely
existence-and-uniqueness.
\end{definition}

Four machine-checked anchor facts frame the landscape:

\begin{itemize}
\item \emph{Rectangles glue over every cover}, uniquely
  (\path{rectangular_cechVanishesCover},
  \path{rectangular_glueCover_existsUnique}).
\item \emph{The trivial cover glues everything}
  (\path{cechVanishesCover_univ}): with one scope containing all
  variables, ``local'' already means ``global.''
\item \emph{The singleton cover recovers the closed theory} of
  Section~\ref{sec:example}
  (\path{cechVanishesCover_singleton_iff}), and the diagonal
  witnesses genuine failure at the cover level
  (\path{diagonal_not_cechVanishesCover}).
\item \emph{The obstruction count refines the condition}
  (\path{coverCechObstruction},
  \path{coverCechObstruction_eq_zero}; the tautological box-gap form
  at the singleton cover is
  \path{cechObstruction_eq_box_sub_encard}).
\end{itemize}

\begin{cohomaside}
The correspondence with the standard apparatus: scopes are the opens
of the ``space'' $V$ (discrete), $\Sec_T$ is a subpresheaf of the
product-of-values presheaf, compatibility is the \v{C}ech
$0$-cocycle condition for the cover, gluing is exactness of $F(V)
\to \Hech^0(U, F)$, and the obstruction count is a cardinality
surrogate for the cokernel.  The linearized versions --- honest
cochain complexes with constant and presheaf coefficients, the
degree-one quotients, and the comparison with
Abramsky--Mansfield--Barbosa's $\gamma$ --- are mechanized and
discussed in Appendix~\ref{app:coefficients}.
\end{cohomaside}

So gluing depends on the cover, monotonely in spirit: coarse covers
glue easily, fine covers glue rarely, and somewhere in between sits a
boundary.  The next section states the machine-checked theorem that
locates the boundary exactly.

\section{Main theorems: gluing works exactly when the pieces fit the
cover}
\label{sec:main}

\subsection{The support--cover gluing theorem}
\label{sec:supportcover}

One more everyday notion.  A constraint $S$ has \emph{support}
$\sigma \subseteq V$ if membership in $S$ depends only on the
coordinates in $\sigma$ (\path{HasSupport}, from \path{Ste.Support})
--- $\sigma$ is the constraint's \emph{scope} in the ordinary CSP
sense.  Say a family $(S_i)$ with scopes $\sigma_i$ \emph{generates}
$T$ if $T = \bigcap_i S_i$, and say its scopes \emph{refine} the
cover $U$ if every $\sigma_i$ fits inside some single scope $U_j$.

In plain language first: \textbf{stitching local answers into global
solutions works, for every compatible family, exactly when the
constraint can be written as a conjunction of pieces each of which
fits inside one scope of the cover.}  If the cover's scopes are big
enough to each swallow whole constraints of some generating family,
gluing is guaranteed; if no generating family fits, some compatible
family is stuck.  There is no third case.

\begin{theorem}[Support--cover gluing; \textbf{mechanized},
\path{Ste.SupportCover}]
\label{thm:supportcover}
Let $T \subseteq \prod_v A_v$ and let $U : J \to \mathcal P(V)$ be a
cover of the variable set.
\begin{enumerate}
\item[(a)] If $T$ is generated by a family whose scopes refine $U$,
  then every compatible family of local sections glues to a global
  solution --- \path{CechVanishesCover}$\,T\,U$ holds
  (\path{cechVanishesCover_of_refiningGenerators}).
\item[(b)] Conversely, if \path{CechVanishesCover}$\,T\,U$ holds,
  then $T = \bigcap_{j} \pi_{U_j}^{-1}\!\bigl(\Sec_T(U_j)\bigr)$: $T$
  is generated by the \emph{cylinder} constraints ``your restriction
  to $U_j$ is a local section'' (\path{sectionCylinder}, supported on
  their scopes by \path{hasSupport_sectionCylinder}) --- a generating
  family whose scopes are exactly the scopes of $U$; indeed gluing is
  \emph{equivalent} to this cylinder identity
  (\path{cechVanishesCover_iff_eq_iInter_sectionCylinder}).
\end{enumerate}
Hence the \v{C}ech obstruction of the cover $U$ vanishes \emph{iff}
$T$ admits a generating family whose scopes refine $U$ --- packaged
as ``gluing $\iff$ there is a per-scope generating family $S_j$ with
\path{HasSupport}$\,(S_j)\,(U_j)$''
(\path{cechVanishesCover_iff_exists_scopedGeneration}).
\end{theorem}

\begin{proof}[Proof idea (exposition only; the machine-checked proof
is in \path{Ste.SupportCover})]
(a) Given a compatible family $(s_j)$, build the only candidate
global assignment there is: for each variable $v$, pick any scope
$U_{j(v)}$ containing $v$ and set $f(v) := s_{j(v)}(v)$.  Agreement
on overlaps makes the choice of scope immaterial, so $f$ restricts to
$s_j$ on \emph{every} scope.  Why is $f$ a solution?  Check the
generators one at a time.  A generator $S_i$ has its whole scope
$\sigma_i$ inside some single $U_{j_0}$; on that scope, $f$ agrees
with $s_{j_0}$, which agrees with the global witness $f_{j_0} \in T$
that makes $s_{j_0}$ a local section.  Since $S_i$ only reads
coordinates in $\sigma_i$, and $f$ matches a known solution there,
$f \in S_i$.  Every generator passes, so $f \in T$.  \emph{The single
place the hypothesis is used is that each generator's scope fits
inside one scope of the cover} --- a constraint straddling two scopes
could read a mixture of two different witnesses and fail.
(Uniqueness of the glue is \path{glueCover_unique}.)

(b) The cylinders visibly contain $T$; conversely an assignment $f$
in all cylinders restricts to a local section on every scope, those
restrictions automatically agree on overlaps (they all come from
$f$), gluing produces a solution $g$ agreeing with $f$ on every
scope, and since the scopes cover $V$, $g = f$ (again
\path{glueCover_unique}), so $f \in T$.
\end{proof}

\begin{remark}[Sanity anchors, all machine-checked]
At the singleton cover, ``a generating family with singleton scopes''
means ``a rectangle,'' and direction (b) specializes to
\path{cechVanishes_iff_rectangular}.  Direction (a) contains
\path{rectangular_cechVanishesCover}: singleton scopes refine
\emph{every} cover.  The diagonal's failure is consistent: its
minimal scope is the full pair $\{x,y\}$
(\path{diagonal_support_full}, in \path{Ste.Support}), contained in
no singleton.  The trivial-cover success
(\path{cechVanishesCover_univ}) is the case where the single scope
swallows all generators.  The theorem also discharges, at the cover
level, the ``residual conjecture'' left open in the docstring of
\path{Ste.VariablePresheaf} (gluing holds precisely for the tractable
fragments) --- provided ``tractable fragment'' is read as
\emph{scope-refinement}, which by Observation~\ref{obs:arity} below
is emphatically \emph{not} computational tractability.
\end{remark}

\begin{remark}[No structural conditions on the cover]
Readers who know the CSP decomposition literature will notice
something missing: classical local-to-global results
(tree-clustering, Vorob'ev-style theorems
\citep{dechterpearl1989tree,freuder1982backtrack}) demand that the
cover's scopes hang together tree-likely (the running intersection
property).  Theorem~\ref{thm:supportcover} imposes nothing of the
sort --- any cover, any overlap pattern.  The reason is the
convention of Section~\ref{sec:fivewords}: STE's local sections are
\emph{globally extendable by definition}.  The classical hypotheses
earn their keep when the local tables are merely \emph{locally}
consistent rather than genuinely extendable
(Section~\ref{sec:pseudo}).
\begin{cohomaside}
In Abramsky--Brandenburger terms: empirical models are arbitrary
compatible presheaf data and can be contextual; the presheaf of
genuine sections of a global set cannot be, in the per-section sense
--- the mechanized \path{amb_extension_always}
(Appendix~\ref{app:twisted}) makes the per-section obstruction
identically zero --- yet family-level gluing still fails exactly when
the cover is too fine for the constraint arity.  This is also why no
nerve-acyclicity hypothesis appears.
\end{cohomaside}
\end{remark}

\subsection{Effective outcome 1: gluing decides representability}
\label{sec:decide}

Theorem~\ref{thm:supportcover} is an \emph{iff}, and each direction
is an effective statement.

\emph{Deciding representability over given scopes.}  ``Can $T$ be
written as a conjunction of constraints living inside the scopes
$U_1, \dots, U_m$?''\ is, by the theorem, the same question as ``does
every compatible family over $U$ glue?''\ --- and direction (b) even
supplies the canonical decomposition when the answer is yes: the
section cylinders (\path{sectionCylinder}), computable from the local
section tables alone.  At the singleton cover this specializes to a
decision procedure for rectangularity
(\path{cechVanishes_iff_rectangular}); at $1$ rectangle it meets the
representation bounds below
(\path{rectCoverNumber_eq_one_iff_cechVanishes}).

\emph{Certifying non-decomposability.}  Contrapositively, a
\emph{single} stuck compatible family over $U$ certifies that no
generating family of $T$ refines $U$ --- no decomposition of the
constraint respects those scopes, so any pipeline treating the scopes
as self-contained is unsound for $T$.  The general contrapositive is
mechanized
(\path{exists_scope_refiningNone_of_not_cechVanishesCover}, in
\path{Ste.SupportCover})\footnote{Mechanized after the first draft of
the underlying exploration memo; the singleton-cover witnesses
(\path{diagonal_not_cechVanishesCover},
\path{allEqual_not_cechVanishes}) instantiate it: no unary
decomposition of these couplings exists.}.

\subsection{Effective outcome 2: where the hardness lives ---
gluing width is arity, not treewidth}
\label{sec:width}

The next two items are \emph{immediate-looking} consequences of the
mechanized theorem, but they are not themselves Lean theorems; they
are the clearly-marked conjectural harvest.

\begin{definition}[Gluing width; not mechanized]
\label{def:gluingwidth}
$\mathrm{gw}(T)$ is the least $w$ such that the sheaf condition
(\path{CechVanishesCover}) holds for $T$ over some cover $U$ all of
whose scopes have at most $w+1$ variables.  By
Theorem~\ref{thm:supportcover}, $\mathrm{gw}(T)$ is exactly (one less
than) the minimal \emph{arity} of a generating constraint family for
$T$: the smallest scopes $T$ decomposes into.
\end{definition}

\begin{observation}[Gluing width is arity, not treewidth --- so
vanishing is not tractability; sketch, not mechanized]
\label{obs:arity}
$\mathrm{gw}$ must not be confused with treewidth, and the library
already contains the separating example.  $\mathrm{allEqual}(n,
\alpha)$ is generated by pairwise-equality constraints --- scopes of
size $2$, so $\mathrm{gw} = 1$ --- while its primal graph is complete
and its induced width is maximal ($n-1$; prose observation in
\path{Ste.CouplingLowerBound}, not a Lean theorem).  More brutally:
graph $3$-colorability is generated by edge-scoped disequality
constraints, so it glues over the edge cover by
Theorem~\ref{thm:supportcover}(a) (this instantiation is prose, not
Lean); deciding $T \ne \emptyset$ is still NP-hard.  \emph{A
vanishing gluing obstruction over the natural cover is compatible
with full computational hardness.}  All the hardness has been pushed
into computing the objects the theorem quantifies over --- the
genuine local sections $\Sec_T(U_j)$ --- and none remains in the
gluing.
\end{observation}

We regard Observation~\ref{obs:arity} as the honest answer to ``is a
vanishing obstruction over a good cover an algorithm?''  It is half
of one: \emph{gluing is the free half; local sections are the
expensive half}.  The algorithmic payoff below is precisely the
theory of when the expensive half is cheap.

\subsection{Effective outcome 3: the algorithmic payoff at bounded
width}
\label{sec:algorithms-sub}

\paragraph{Bag tables are local sections.}  The treewidth chain
(\path{Ste.Treewidth} through \path{Ste.JunctionTree}) is built on
the \path{table} of a constraint on a scope $\sigma$: the set of
restrictions to $\sigma$ of solutions.  Comparing definitions,
\path{table}~$\sigma$~$T$ and \path{localSections}~$T$~$\sigma$ are
the \emph{same construction} --- both are the image of $T$ under
restriction --- so the machine-checked junction-tree theorems are,
without change, statements about local sections over the bag
cover:\footnote{The one-line identification lemma is dispatched
mechanization work (Section~\ref{sec:conclusion}); the two
definitions are syntactically the same image construction.}

\begin{itemize}
\item \emph{Faithfulness: the bag tables jointly remember the
  constraint.}  A scoped constraint is exactly the preimage of its
  bag table (\path{HasSupport.eq_preimage_table}), and a global
  assignment is a solution iff it restricts into every bag table
  (\path{mem_joinConstraint_iff_restrict},
  \path{junctionTables_faithful}).  This is
  Theorem~\ref{thm:supportcover}(b) in concrete clothing: membership
  in $T$ is decided scope-by-scope.
\item \emph{Cost: local sections are small when scopes are small.}  A
  bag of at most $w+1$ variables over alphabets of size at most $k$
  tables in at most $k^{w+1}$ rows (\path{table_encard_le_pow}); a
  width-$w$ elimination order materializes at most $n \cdot k^{w+1}$
  rows in total (\path{bucketEliminate_total_space},
  \path{junctionTree_size_le}; linear-in-$n$ form
  \path{junctionTree_size_linear} at the minimal duplicate-free
  width \path{junctionWidth}).
\item \emph{Computability: bucket elimination computes genuine
  sections at bounded width.}  Projective bucket elimination computes
  exactly the projection of the solution set onto the remaining
  variables at each step
  (\path{joinConstraint_projectBucketEliminate} and neighbours in
  \path{Ste.BucketConsistency}) --- genuine local sections, not
  approximations --- decides feasibility
  (\path{projectBucketEliminate_decides},
  \path{bucketEliminate_decides}), and extracts a solution
  backtrack-free
  (\path{exists_extension_of_mem_projectBucketEliminate}); the fused
  statement is \path{projectBucketEliminate_complete_solver} and, for
  adaptive consistency proper, \path{decreasingRun_adaptiveSolver}
  \citep{dechter1999bucket,dechter2003constraint,dechterpearl1987network}.
\end{itemize}

So the mechanized algorithmic layer already computes the genuine
local-section tables over an elimination cover, in $O(n \cdot
k^{w+1})$ at induced width $w$ (\path{inducedTreewidth},
\path{achievesWidth_inducedTreewidth}; the two-sided bridge to
Robertson--Seymour treewidth \citep{robertson1986graphminorsII} is
\path{treewidth_primalGraph_le} and
\path{inducedTreewidth_le_treewidth_sub_one}; nonserial
dynamic-programming ancestry in \citealp{bertele1972nonserial}).
Combined with Theorem~\ref{thm:supportcover}(a) --- gluing is free
once the generators fit the bags --- this is a complete
local-to-global pipeline for bounded-width instances: $O(n \cdot
k^{w+1})$ preprocessing replaces $2^N$ enumeration, and bag-local
queries read off the section tables soundly by faithfulness.

\paragraph{Consistency algorithms compute pseudo-sections.}
\label{sec:pseudo}
Local-consistency algorithms
\citep{mackworth1977consistency,freuder1982backtrack} compute
something weaker than genuine local sections: tables closed under
local checks, with no guarantee their entries extend to global
solutions.  Call these \emph{pseudo-sections}.  The gap between
pseudo and genuine is exactly where the hardness displaced by
Observation~\ref{obs:arity} ends up living, and the mechanized
tractability results are precisely theorems that \emph{pseudo equals
genuine} under structural hypotheses: on a chain, arc consistency
puts every surviving value on a global solution
(\path{arcConsistent_backtrackFree},
\path{arcConsistent_chain_solvable}); on a rooted tree, directional
arc consistency does the same (\path{treeArcConsistent_solvable},
\path{Ste.ConsistencyTree}); at induced width $w$, directional
consistency of the buckets does the same, and bucket elimination
establishes it while preserving solutions
(\path{directionalConsistent_solvable} in
\path{Ste.AdaptiveConsistency}; \path{Ste.BucketConsistency}).

\begin{center}
\footnotesize
\begin{tabular}{@{}p{9em}p{13em}p{14em}@{}}
\toprule
 & pseudo-sections ($k$-consistent tables) & genuine sections
 (\path{localSections}) \\
\midrule
cheap to compute & yes (local closure) & only at bounded width \\
entries globally extendable & not guaranteed & always, by definition
 (cf.\ \path{amb_extension_always}) \\
family gluing & fails generically & iff generators fit the cover
 (Thm.~\ref{thm:supportcover}) \\
where hardness sits & pseudo $\ne$ genuine gap & computing
 $\Sec_T$ \\
\bottomrule
\end{tabular}
\end{center}

\noindent The conjectural summary (a reading, not a theorem):
\emph{the cohomology-shaped content of STE hardness is the
pseudo/genuine gap, and every tractability theorem in the library is
a statement that the gap closes over a structured cover.}

\subsection{Effective outcome 4: representation bounds and
certificates}
\label{sec:quant-sub}

A tempting slogan --- ``the bigger the obstruction, the bigger the
forced representation blow-up'' --- is settled \emph{negatively}, and
machine-checked, in \path{Ste.RepresentationBounds}.  The measuring
stick is the \emph{rectangle-cover number} $\rectrho(T)$
(\path{rectCoverNumber}): the least number of no-coupling
(rectangular) constraints whose \emph{union} is exactly $T$ --- ``how
many coupling-free alternatives describe $T$?''  This is the STE
analogue of the nondeterministic communication-complexity cover
number \citep{kushilevitz1997communication,aho1983notions}, with the
LP-extension-complexity lineage
\citep{yannakakis1991expressing,fiorini2015exponential} as outlook.
Machine-checked:

\begin{itemize}
\item \emph{Bounds.}  A \emph{fooling set} --- solutions no two of
  which can be mixed coordinate-wise and stay solutions --- forces
  one rectangle each: $\mathrm{fool}(T) \le \rectrho(T) \le |T|$
  (\path{foolingNumber_le_rectCoverNumber},
  \path{IsFoolingSet.encard_le_of_hasRectCover},
  \path{rectCoverNumber_le_encard}); the mixing engine is
  \path{IsRectangle.mix_mem}.
\item \emph{The bottom of the scale is the gluing condition}:
  $\rectrho(T) = 1 \iff$ the singleton-cover obstruction vanishes
  (\path{rectCoverNumber_eq_one_iff_cechVanishes},
  \path{one_lt_rectCoverNumber_iff}).
\item \emph{Exact values on the coupling}:
  $\rectrho(\mathrm{allEqual}(n,\alpha)) = |\alpha|$
  (\path{rectCoverNumber_allEqual}; fooling set
  \path{isFoolingSet_allEqual}, \path{foolingNumber_allEqual}).
\item \emph{No law connects obstruction size to $\rectrho$}
  (\path{not_forall_cechObstruction_le_rectCoverNumber},
  \path{not_forall_rectCoverNumber_le_cechObstruction}): rectangles
  give $\obs = 0 < 1 = \rectrho$; $\mathrm{allEqual}(3,
  \mathrm{Bool})$ gives $\rectrho = 2 < 6 = \obs$
  (\path{rectCoverNumber_lt_cechObstruction_allEqual}); on the
  two-variable diagonal the two coincide
  (\path{diagonal_cechObstruction_eq_rectCoverNumber}) --- a
  coincidence of the smallest example, not a law.
\end{itemize}

Two representation grammars are in play: $\rectrho$ measures
$\cup$-of-boxes (DNF-like) representations; the junction-tree
machinery measures $\cap$-of-cylinders (CNF-over-bags)
representations.  The machine-checked triple $(\obs, \rectrho,
\text{width})$ on $\mathrm{allEqual}$ shows all three diverge:
obstruction exponential, $\rectrho$ constant, width maximal.

\begin{conjecture}[$\rectrho$ vs.\ width incomparability, other
direction; not mechanized]
A width-$1$ chain of disequality constraints over an alphabet of
size $k \ge 3$ on $n$ variables has feasible set of size
$k(k-1)^{n-1}$ and, we conjecture, $\rectrho = 2^{\Omega(n)}$ (a
fooling-set construction along alternating assignments looks
plausible).  If so, $\rectrho$ and treewidth are incomparable in both
directions: $\rectrho$ answers ``how many boxes,'' width answers
``how large a bag.''
\end{conjecture}

\emph{Practitioner's reading}: fooling sets are cheap-to-exhibit
certificates that no small union-of-boxes representation exists;
stuck compatible families over a proposed cover
(Section~\ref{sec:decide}) are cheap-to-exhibit certificates that no
generating family respects those bags.  Both certificate kinds are
mechanized at the definition level and computed exactly on the
coupling family.  Cheap NO-certificates steer the expensive
machinery: test a candidate bag cover with a small stuck-family
search before paying for elimination.

\section{Carlson connection points}
\label{sec:carlsonconn}

Does Carlson's cryptanalytic STE \citep{carlson2012} map onto the
sheaf picture?  The honest answer is: \emph{partially, and the
partial map is instructive}.  We separate three tiers.

\subsection{What maps, with mechanized witnesses}

\emph{Observation accumulation is the constraint axis.}  Carlson's
``more intercepted traffic narrows the key search''
(\path{feasibleKeys_antitone}) is an instance of the axis that glues
for free (Section~\ref{sec:why}, item 1;
\path{partialFeasibilitySet_antitone},
\path{feasibilitySet_eq_iInter_cover}): observations index property
sets, and covers of the \emph{observation} index can never obstruct.
The dissertation's accumulation arguments live entirely on this
axis, which is why they need no gluing hypotheses.

\emph{Residual key counting is fiber counting.}  Carlson's Lemma~4.1
(\path{card_consistent_keys}) counts the permutations agreeing with
the true key $\sigma$ on the observed symbol set $T$: $(|A|-|T|)!$.
In the paper's vocabulary, a key's action on $T$ is a partial
assignment on the scope $T$ (variables $=$ alphabet symbols, values
$=$ alphabet symbols), and the count is the number of \emph{global
witnesses of one local section} --- the fiber of the restriction map
over $\sigma|_T$.  Unicity (\path{Unicity}) is the statement that
accumulated observations shrink every such fiber to one point.  This
reading is faithful to the mechanized statement as it stands; its
formal restatement through the restriction map used by
\path{localSections} is now itself mechanized on a research branch
(Section~\ref{sec:scopedobs}).

\subsection{The scoped-observation theorems (mechanized on a research
branch)}
\label{sec:scopedobs}

The genuinely sheaf-shaped content of the substitution cipher is
this: \emph{an observed ciphertext constrains the key only on the
symbols that occur in it} --- an observation constraint is a
\emph{scoped} constraint in the sense of \path{HasSupport}, with
scope the ciphertext's symbol set (on the decryption-side
parametrization of keys).  Meanwhile the requirement that the key be
a \emph{permutation} --- injectivity across all symbols --- is a
global \emph{coupling}, the cryptanalytic cousin of
$\mathrm{allEqual}$'s all-different flavor, and it is provably
non-rectangular.  So Carlson's instance factors exactly along this
paper's main axis: scoped evidence (gluable, by
Theorem~\ref{thm:supportcover}(a), over any cover the symbol sets
refine) against one global coupling (the permutation constraint,
which no proper cover swallows).

These statements were dispatched during the writing of this paper
and have been \textbf{mechanized, sorry-free, on the research branch
\texttt{research/carlson-sheaf-bridge}} (commit \texttt{4ae7425},
module \texttt{Ste.Carlson.SheafBridge}; independently verified to
contain no \texttt{sorry} and to compile against the project ---
axioms limited to the standard three).  The branch is \emph{not yet
merged} into the main library, so per the bright line these four
results are cited with the branch stated and are kept out of the
main-library tally: \texttt{hasSupport\_decConstraint} (the
decryption-side observation constraint \texttt{decConstraint} has
support exactly the ciphertext's symbol set);
\texttt{mem\_keyPropertySet\_iff\_decConstraint} (it is Carlson's
\path{keyPropertySet} of the substitution \texttt{subCipher}, read
through key inversion); \texttt{bijKeys\_not\_rectangular} (the
permutation coupling \texttt{bijKeys} is not rectangular over any
nontrivial alphabet); and \texttt{card\_restrict\_fiber} (Carlson's
Lemma~4.1 restated through the restriction map used by
\path{localSections}: the permutations restricting on $T$ to a given
local section number $(|A|-|T|)!$, derived from
\path{card_consistent_keys}).  With these, the fiber-counting reading
of Section~\ref{sec:carlsonconn} is a theorem, not an analogy ---
modulo the merge.

\subsection{What does not map}

Two honest negatives.  First, Carlson's hypothesis space is a key
space, not a product of variables; the sheaf reading above requires
re-embedding keys into the function space $A \to A$ and carrying the
permutation constraint separately.  Nothing in the dissertation uses
cover-locality on the key's coordinates; the productive structure
there is \emph{class-level reduction} --- embedding one cipher
class's keys in another's
(\path{permutation_reduces_to_substitution},
\path{coperm_injective},
\path{substitution_not_reducible_to_permutation},
\path{psp_reduces_to_substitution}) --- which is group theory, not
gluing, and we see no gluing-theoretic content in it.  Second, the
dissertation's information-theoretic layer (unicity distances,
entropy of meaningful text) has no mechanized counterpart in this
project and no sheaf reading we would defend; connecting the
obstruction counts of Section~\ref{sec:quant-sub} to
unicity-distance estimates is \emph{speculation} and is recorded as
such in Section~\ref{sec:conclusion}.

\section{Empirical analysis and constructive proof-of-work}
\label{sec:empirical}

The project's empirical layer runs a \emph{verified} STE solver on
frame corpora extracted from historical texts.  This section reports
what ran, exactly what the solver computes, and what the
cohomological theory adds --- with the pipeline's trust boundary
stated plainly.

\subsection{The frame-corpus constraint class and the verified
solver}

An empirical corpus is a \emph{frame} family: each author $a$ fixes a
partial assignment $\mathrm{frame}\,a : V \to \mathrm{Option}\,W$ of
values to variables (silence $=$ \textsf{none}), and the property set
of $a$ is the set of total assignments agreeing with $a$ wherever $a$
speaks (\path{frameConstraint}, in \path{Ste.FiniteInstance}).  For
this class the feasibility question collapses to a decidable pairwise
check, and the collapse is machine-checked: the corpus is jointly
satisfiable iff no two authors assign the same variable different
values (\path{Consistent}, with the bridge
\path{nonempty_iff_consistent} and the verdict form
\path{feasibilitySet_eq_empty_iff}); the quantitative refinement
counts distinct values per variable (\path{disagreementDegree},
characterized by \path{agree_iff_degree_le_one}).

The executable (\path{lean/TidesExe.lean}) parses a
\texttt{frames.json} at runtime and evaluates exactly these verified
definitions --- \path{STE.Consistent} and
\path{STE.disagreementDegree} --- on the parsed frames.  \textbf{To
be precise about what ran: the solver's algorithm is the pairwise
consistency check of \path{Ste.FiniteInstance}; it imports no
cohomology module, and no \v{C}ech computation executed on any
corpus.}  The trust boundary is the JSON parser and Lean's evaluator;
the aggregation logic is the verified one.

\subsection{The runs}

Three corpora, one pipeline (LLM extraction of frames from period
texts; the verdicts computed by the verified layer):

\begin{itemize}
\item \emph{Tides} (\path{Ste.TidesCorpus},
  \path{Ste.TidesComputable}): pre-Newtonian authors on the cause of
  the tides, canonicalized to a $7$-variable schema.  Verdict:
  globally \emph{inconsistent}, and pervasively so --- the
  five-author ``attraction'' camp is internally consistent on the
  moon's role but not overall
  (\path{tides_inconsistent},
  \path{attraction_camp_moon_consistent}); the kernel-computed
  verdict (\path{feasible_eq_empty_computed}) agrees with the
  hand-proved one, and the corpus is definitionally an instance of
  the general frame class (\path{constraint_eq_frameConstraint}).
\item \emph{Federalist} (86 essays; \citealp{hamilton1788federalist}):
  $10$ variables; inconsistent but \emph{narrowly} --- $7$ of $10$
  variables unanimous where spoken, exactly $3$ conflict sites
  carrying the whole obstruction.
\item \emph{Ratification, two camps} (the $86$ Federalist essays plus
  $73$ Anti-Federalist essays; $159$ documents): the consistent core
  shatters --- all $10$ variables become conflict sites.  The
  contrast $3 \to 10$ is the sharp empirical signature of adding an
  opposing camp.
\end{itemize}

\subsection{What the cohomology adds: why the cheap check is the
right one}

Here is the honest division of labor.  The solver ran a
\emph{pairwise} check.  The theory of this paper explains \emph{why
pairwise is enough for this constraint class} --- why no
higher-order gluing failure can hide in a frame corpus.

The plain argument: each author's \path{frameConstraint} is, by
construction, a conjunction of \emph{unary} constraints (``variable
$v$ has value $x$,'' one per spoken variable --- read the definition:
$\forall v\,x$, spoken $\Rightarrow$ agreement).  Unary scopes refine
\emph{every} cover, so by the mechanized
Theorem~\ref{thm:supportcover}(a)
(\path{cechVanishesCover_of_refiningGenerators}) the feasible set of
a frame corpus glues over every cover of the variables; equivalently,
the corpus feasible set is a rectangle, the class for which
per-variable and pairwise reasoning is complete
(\path{cechVanishes_iff_rectangular},
\path{rectangular_cechVanishesCover}).  There are no phantoms to
fear: locally consistent readings of a frame corpus are always
globally realizable, which is exactly why the decidable surrogate
\path{Consistent} can coincide with genuine feasibility
(\path{nonempty_iff_consistent} --- itself machine-checked,
independent of any cohomology).  The \emph{formal} packaging of the
rectangularity argument (``\path{frameConstraint} is rectangular,
hence the corpus feasible set is, hence it \v{C}ech-vanishes for
every cover'') is exact, Lean-ready, and dispatched
(Obligation~\ref{ob:frame}); until it lands, the machine-checked
pillars are the two endpoints already cited ---
\path{nonempty_iff_consistent} on the solver side and
Theorem~\ref{thm:supportcover} on the theory side --- and the
connecting rectangularity lemma is the marked gap.

Two consequences worth stating plainly.  First, the empirical runs
are a genuine, if structurally easy, instance of the paper's subject:
a local-reasoning solver whose completeness is exactly a gluing fact.
Second, the interesting failure mode --- coupled constraints, where
pairwise reasoning is \emph{not} complete
(Section~\ref{sec:example}) --- does not arise for frame corpora, but
will arise the moment extraction produces \emph{relational} claims
(``$v_1$ and $v_2$ stand in relation $R$'') rather than
variable-value assignments.  A solver for that regime, carrying its
gluing certificate explicitly (compute bag-local sections, check
compatibility, invoke Theorem~\ref{thm:supportcover} for
soundness), is the natural constructive next step; \textbf{no such
cohomology-labelled solver exists yet}, and we flag it as future
work rather than claim it.

\section{Conclusion and future work}
\label{sec:conclusion}

\subsection{What is machine-checked, in one paragraph}

The constraint side of STE glues for free
(\path{feasibilitySet_eq_iInter_cover}); the corpus side glues with
an exact computable correction (\path{mem_feasible_union_iff}); on
the variable side, per-variable gluing characterizes the absence of
coupling exactly (\path{cechVanishes_iff_rectangular}),
coupling-free constraints glue over every cover with unique glues
(\path{rectangular_glueCover_existsUnique}), couplings genuinely fail
(\path{diagonal_not_cechVanishesCover}) with exactly computed,
exponentially growing phantom counts
(\path{cechObstruction_allEqual},
\path{cechObstruction_allEqual_ge}), and the boundary is exact: the
support--cover gluing theorem (Theorem~\ref{thm:supportcover},
\path{Ste.SupportCover}) with its certificate contrapositive
(\path{exists_scope_refiningNone_of_not_cechVanishesCover}).  Cover-shape-only invariants
are provably blind (\path{cechH1_subsingleton},
Appendix~\ref{app:constant}) while the section-aware linear
obstruction fires over every nontrivial ring
(\path{mixedFamily_not_linearlyGlues}) even though the per-section
AMB obstruction cannot (\path{amb_extension_always}).  The
rectangle-cover number is sandwiched, characterized at $1$ by
gluing, computed on the coupling, and provably incomparable to the
obstruction magnitude
(\path{not_forall_cechObstruction_le_rectCoverNumber} and mirror).
The bounded-width stack --- tables, bucket elimination, junction
trees, adaptive consistency --- is mechanized with faithfulness,
$O(n \cdot k^{w+1})$ size, and a fused complete-solver statement
(\path{projectBucketEliminate_complete_solver}).  Carlson's STE
recast, counting lemma, and reduction chain are mechanized
(\path{Ste.Carlson}); the frame-corpus solver equivalence is
mechanized (\path{nonempty_iff_consistent}) and ran on three corpora.

\subsection{What remains conjectural}

The gluing-width invariant and its arity reading
(Definition~\ref{def:gluingwidth}, Observation~\ref{obs:arity}); the
pseudo/genuine-gap reading of tractability
(Section~\ref{sec:pseudo}); the $\rectrho$-vs-width incomparability
conjecture; the bag-local linear relaxation with exactness controlled
by the twisted invariant (Appendix~\ref{app:twisted}); and any
connection between obstruction counts and unicity distance
(Section~\ref{sec:carlsonconn}) --- explicitly speculation.

\subsection{Dispatched proof obligations}

The following exact obligations are, at the time of writing,
dispatched to in-flight Lean mechanization efforts on named research
branches.  \emph{None of them is cited as proven anywhere in this
paper}; when they land sorry-free, the noted sections upgrade.

\begin{obligation}[Carlson sheaf bridge; branch
\texttt{research/carlson-sheaf-bridge} --- \textbf{discharged},
awaiting merge]
\label{ob:carlson}
(a) The decryption-side observation constraint of the substitution
cipher has \path{HasSupport} equal to the symbol set of the
ciphertext.  (b) Its equivalence with \path{keyPropertySet} of the
substitution \path{CipherSystem} via key inversion.  (c) The
permutation coupling $\{g : A \to A \mid g \text{ bijective}\}$ is
not rectangular over a nontrivial alphabet.  (d) Carlson's Lemma~4.1
restated through the restriction map (from
\path{card_consistent_keys}).  \emph{Status: all four items
mechanized sorry-free on the branch at commit \texttt{4ae7425}
(Section~\ref{sec:scopedobs}); they enter the main-library tally on
merge.}
\end{obligation}

\begin{obligation}[Frame rectangularity; branch
\texttt{research/frame-rectangular}]
\label{ob:frame}
(a) \path{IsRectangle} of each author's \path{frameConstraint}; (b)
\path{IsRectangle} of the corpus feasible set
$\mathrm{feasibilitySet}(\mathrm{frameConstraint})$; (c) hence
\path{CechVanishesCover} for every cover of the variables; (d)
optionally, the unary-generation form feeding
\path{cechVanishesCover_of_refiningGenerators}.  Upgrades
Section~\ref{sec:empirical} from prose argument to cited theorem.
\end{obligation}

\begin{obligation}[Near-term backlog; branch
\texttt{research/mechanize-conjectures}, in flight independently of
this paper]
\label{ob:backlog}
The \path{table} $=$ \path{localSections} identification; a
pair-cover gluing statement for $\mathrm{allEqual}$ and a Lean-level
gluing-width definition; and the second-tier representation bounds.
Upgrades Sections~\ref{sec:width} and~\ref{sec:algorithms-sub}.
\end{obligation}

\subsection{Further directions}

A pseudo-section presheaf as a first-class object, with the unifying
vanishing theorem ``running-intersection cover $+$ directional
consistency $\Rightarrow$ pseudo $=$ genuine'' subsuming the chain,
tree, and adaptive cases; the variable-side analogue of the corpus
axis' exact correction term (splice two partial solutions, collect
straddling constraints), then its $k$-piece iteration; a bag-local
must/may/cannot query theorem over junction tables; the twisted
degree-one theory (Appendix~\ref{app:twisted}) toward exactness
guarantees for linear relaxations in the style of
\citet{abramsky2017contextualfraction}; and the relational-claims
solver of Section~\ref{sec:empirical} --- the first solver that
would carry its gluing certificate explicitly.  Pleasingly for a
paper about cohomology, the content of its main theorem is that in
set-theoretic estimation the cohomology, properly aimed, is exactly
the bookkeeping of who is allowed to talk to whom --- and that
bookkeeping is machine-checked.

\appendix

\section{For the cohomologist: what the coefficients see}
\label{app:coefficients}

This appendix is the paper's heavy-machinery annex.  Nothing in the
main text depends on it beyond the two plain-language morals already
stated there: cover-shape-only invariants are provably blind
(\ref{app:constant}), and a section-aware linear invariant fires one
degree lower than the AMB program would suggest (\ref{app:twisted}).
The Abramsky--Brandenburger program grades gluing failure
cohomologically by linearizing the section presheaf
\citep{abramsky2011sheaf,abramsky2012cohomology}; the library
mechanizes two coefficient regimes, and the contrast between them is,
in our view, the single most instructive machine-checked fact in
this layer.

\subsection{Constant coefficients see nothing --- a machine-checked
no-free-lunch}
\label{app:constant}

\path{Ste.CechComplex} builds the constant-coefficient cochain
complex on the full nerve of a cover index --- modules $C^0 = \iota
\to M$, $C^1 = \iota \to \iota \to M$, \dots\ with the standard
alternating coboundaries (\path{cechD0}, \path{cechD1},
\path{cechD2}), the complex identities (\path{cechComplex_d_comp_d},
\path{cechComplex_d2_comp_d1}), and honest quotient modules
$\Hech^1, \Hech^2$ (\path{cechH1}, \path{cechH2}).  The vanishing
theorem (\path{cechH1_subsingleton}, \path{cechH2_subsingleton})
proves $\Hech^1 = \Hech^2 = 0$ over any nonempty index type: the
full nerve is a simplex, and a simplex is acyclic; $\Hech^0$ is
exactly the coefficients (\path{cechH0EquivCoeff}).

The computational moral: \emph{any invariant computed from the
combinatorics of the cover alone --- ignoring the section data ---
is useless for detecting coupling}; the machine checks that it is
identically zero.  Whatever a ``cohomological speedup'' means, its
input must include the tables of local sections --- exactly the data
the CSP algorithms of Section~\ref{sec:algorithms-sub} already
manipulate.  This is consistent with the contextuality literature,
where the interesting classes live in relative cohomology with
\emph{presheaf} coefficients
\citep{abramsky2012cohomology,caru2017cohomology}.

\subsection{Twisted coefficients see the coupling --- one degree
lower than expected}
\label{app:twisted}

\path{Ste.TwistedCech} builds the honest twisted complex:
coefficients in the linearized section presheaf $F_R(W) =
(\Sec_T(W)) \to_{\!f} R$ (free $R$-module on genuine local sections,
\path{twistedCoeff}), restriction by pushforward along section
restriction (\path{sectionRes}, \path{twistedRes}, functorial by
\path{twistedRes_twistedRes}), coboundaries \path{twistedD0},
\path{twistedD1} with $d^1 \circ d^0 = 0$
(\path{twisted_d1_comp_d0}), and $\Hech^0$ as \path{twistedH0}.  Two
machine-checked results locate the STE coupling precisely.

\begin{enumerate}
\item \emph{The AMB per-section obstruction degenerates in STE}
  (\path{amb_extension_always}).  Abramsky--Mansfield--Barbosa
  attach to a single local section $s_1$ a class $\gamma(s_1) \in
  \Hech^1$ vanishing iff $s_1$ extends to a compatible family of the
  linearized presheaf \citep[Prop.~4.2]{abramsky2012cohomology}.
  The mechanization proves that in STE this extension condition
  holds for \emph{every} section of \emph{every} constraint over
  \emph{every} cover --- because \path{localSections} only ever
  contains restrictions of global solutions, the extension witness
  is free.  The $\Hech^1$-graded per-section obstruction is
  structurally blind here: STE's section presheaf bakes global
  extendability into its stalks.  (The long exact sequence and the
  literal connecting map $\gamma$ are \emph{not} mechanized; what is
  proven is the extension condition AMB prove equivalent to $\gamma
  = 0$.)
\item \emph{The obstruction that fires lives in the cokernel of
  $F_R(V) \to \Hech^0(U, F_R)$}
  (\path{mixedFamily_not_linearlyGlues},
  \path{diagonal_mixed_cocycle_not_in_globalRange}).  A compatible
  family's point-mass cochain is a twisted $0$-cocycle
  (\path{familyCochain_mem_twistedH0}); the family \emph{glues
  $R$-linearly} (\path{LinearlyGlues}) iff that cocycle lifts to a
  formal $R$-combination of global solutions
  (\path{linearlyGlues_iff_mem_range}).  For the two-bit coupling
  and its singleton cover, the mixed family fails to glue linearly
  over \emph{every nontrivial commutative ring} --- negative and
  rational coefficients included.  This is strictly stronger than
  the set-level failure (\path{diagonal_gluing_fails}): in the AMB
  world every no-signalling model has a global section over the
  signed reals \citep{abramsky2011sheaf}, but the STE coupling is
  not even signed-linearly explainable by global data.  Set-level
  gluing always implies linear gluing
  (\path{GluesCover.linearlyGlues}), and rectangles glue linearly
  over every cover (\path{rectangular_linearlyGlues}), so
  nonmembership is a genuine obstruction invariant
  (\path{diagonal_twisted_obstruction_detects}).
\end{enumerate}

Whether the twisted $\Hech^1$ (carried as the submodule pair
\path{twistedCoboundaries1} $\le \ker$ \path{twistedD1}, with
triviality expressed as \path{TwistedH1Trivial}; the literal
quotient type is blocked by a typeclass diamond, documented in the
file) vanishes for pairwise-disjoint covers is stated as \emph{open}
in the source and remains open here.  It is the first stepping stone
toward exactness guarantees for bag-local linear relaxations ---
signed/affine analogues of the $k$-consistency hierarchy, cousins of
the contextual-fraction linear programs of
\citet{abramsky2017contextualfraction} --- with the standing caution
from the literature
\citep{caru2017cohomology,okay2017topological,aasnaess2022comparing}
that invariants of this style can have false negatives; the
mechanized \path{amb_extension_always} is precisely such a
false-negative theorem.

\begin{cohomaside}
On the corpus axis the presheaf is separated (identical ambient
hypothesis on overlaps), so its theory closes at $\Hech^0$; a
Mayer--Vietoris organization of iterated corpus splits (standard
machinery, e.g.\ \citealp{bott1982forms}) would carry correction
terms that are precisely the \path{crossObstruction} sets, and
separatedness makes us expect immediate degeneration --- a polite
way of saying the corpus axis needs no cohomology beyond what is
mechanized.  A bi-indexed presheaf on $\text{(corpora)} \times
\text{(scopes)}$, corpus direction separated, variable direction
carrying the twisted theory, is (frankly speculative) outlook.
\end{cohomaside}

\bibliographystyle{plainnat}
\bibliography{../../refs}

\end{document}
